\documentclass[10pt,a4paper]{amsart}
\usepackage[margin=19mm]{geometry}
\usepackage{amsmath,amssymb,mathtools}
\usepackage[scr=rsfs,cal=euler]{mathalfa}
\usepackage{xfrac}
\usepackage[unicode,hidelinks,bookmarks=false]{hyperref}
\newcommand{\ie}{i.e.\ }
\newcommand{\cf}{cf.\ }
\newcommand{\resp}{respectively\ }
\newcommand{\Subsection}{\subsection}
\providecommand{\nobreakdash}{\nobreak\hbox{-}}
\setlength{\emergencystretch}{2em}
\multlinegap0pt
\usepackage{bm}
\usepackage{bbm}
\usepackage{dsfont}
\usepackage{xspace}
\usepackage{cancel}
\usepackage{mathtools}
\usepackage{amsthm}
\makeatletter
\@ifundefined{thm}{\newtheorem{thm}{Theorem}}{}
\makeatother
\def \e {{\bf e}}
\newcommand{\ob}[1]{\left(#1\right)}
\title[Quantum walks on finite and bounded infinite graphs]{Quantum walks on finite and bounded infinite graphs}
\author{Chris Godsil, Steve Kirkland, Sarojini Mohapatra, Hermie Monterde,, Hiranmoy Pal}
\date{Source: arXiv:2510.05306v1 (CC BY 4.0)}
\begin{document}
\maketitle

\noindent\emph{Source and licence.} Quantum walks on finite and bounded infinite graphs. Version arXiv:2510.05306v1 (CC BY 4.0), distributed under the Creative Commons Attribution 4.0 International licence, as recorded in the arXiv licence metadata for this exact version and shown on the abstract page https://arxiv.org/abs/2510.05306v1. Full rights notice: licenses/arxiv-2510.05306-CC-BY-4.0.txt.\par\smallskip

\noindent\textit{Editorial scope.} Counted unit: the Thm whose source label is \texttt{godgen} in the article (source0.tex lines 321--339, source label \texttt{godgen}), delivered together with the article's own complete original proof of it (source0.tex lines 340--389, one \texttt{proof} environment of 50 lines). No mathematics is written, repaired or re-derived: statement and proof are the article's own text line for line. Required context retained uncounted: none. Every definition, display and label of those lines is retained unchanged. Omitted from this package and not redistributed: 1--320, 390--1391. Nothing inside a retained excerpt is omitted or altered: every retained body is delivered exactly as the article's lines are written, so no omission receipt of any kind accompanies this package. Citations: the retained excerpts carry no citation command at all, so no bibliography entry of the article is transcribed and no reference is redistributed. Reference adaptation: none was needed and none was made; every reference inside the delivered unit resolves inside it, so no unresolved reference or citation is produced. Printed slips kept literal: no printed word, symbol, hypothesis or number of the counted statement or of its proof was altered. Layout and class: the article's own class is \texttt{article} with packages including etex, inputenc, amsmath, amssymb, amsthm, amsfonts, mathrsfs, bbm, commath, multicol, blindtext, graphicx, tikz-network, subcaption; the wrapper is the lane's pinned \texttt{amsart} editorial wrapper at 10pt on A4, which restates the theorem-like environments the retained text needs and adds the article's own macro definitions that the retained text needs (\texttt{\textbackslash e}, \texttt{\textbackslash ob}), with extra wrapper packages (bm, bbm, dsfont, xspace, cancel, mathtools). The delivered numbering is the unit's own. Assets: no retained span contains a figure environment, a graphics-inclusion command, a PDF-page inclusion, a table or a TikZ picture; no image file is copied into this package, the package declares no asset dependency and no image file is redistributed with it. Licence: version arXiv:2510.05306v1 of this article is distributed under the Creative Commons Attribution 4.0 International licence, as recorded in the arXiv licence metadata for this exact version and shown on the abstract page https://arxiv.org/abs/2510.05306v1. A full rights notice is delivered at licenses/arxiv-2510.05306-CC-BY-4.0.txt. Characters: every retained excerpt is pure ASCII, so the delivered engine, XeLaTeX with its default Latin Modern text font, reports no missing character.
\begin{thm}
\label{godgen}
Let $G$ be a bounded graph with edge-perturbed twin subgraphs $X_1$ and $X_2$ with $A=A(G)$, and $\Pi$ be a partition of $G$ where $\{a,\hat{f}(a)\}$ is a cell for each $a\in V (X_1)$. Let $Q=[\mathcal{B}~~\mathcal{C}]$ be the matrix where the columns of $\mathcal{B}$ consist of all vectors from the set $\{\frac{1}{\sqrt2}(\e_a-\e_{\hat{f}(a)}):a\in V(X_1)\}$ and $\mathcal{C}$ is the normalized characteristic matrix of $\Pi.$ The following are equivalent:
\begin{enumerate}
\item $\Pi$ is equitable.
\item The column space of $Q$ is $A$-invariant.
\item There is a matrix $B$ such that $AQ=QB$.
\item The matrices $A$ and $QQ^T=\displaystyle I_{2|V(X_1)|}\oplus \left(\bigoplus_{k>|V(X_1)|} \frac{1}{|V_k|}J_{|V_k|}\right)$ commute.
\end{enumerate}
Moreover, if such a matrix $B$ exists, then
\begin{equation}
\label{B}
B:=\begin{bmatrix}
     A(X_1)-A' & \mathbf{0}\\
     \mathbf{0} & A(G/\Pi)
     \end{bmatrix},
\end{equation}
and $A'$ is a symmetric matrix whose rows and columns are indexed by the vertices of $X_1$ and $X_2,$ respectively, and $(A')_{u,v}\neq 0$ if and only if $u$ and $v$ are adjacent in $G$.
\end{thm}

\begin{proof}
We first prove $1\Leftrightarrow 3$. Assume $\Pi$ is equitable so that $A\mathcal{C}=\mathcal{C}A(G/\Pi)$. As $X_1$ and $X_2$ are edge-perturbed twin subgraphs in $G$ and $\{a,\hat{f}(a)\}$ is a cell for all $a\in V (X_1)$, we get $A\mathcal{B}=\mathcal{B}(A(X_1)-A')$. Combining the two preceding equations gives us
\begin{equation*}
AQ=[A\mathcal{B}~~A\mathcal{C}]=[\mathcal{B}(A(X_1)-A')~~\mathcal{C}A(G/\Pi)]=[\mathcal{B}~~\mathcal{C}]\begin{bmatrix}
     A(X_1)-A' & \mathbf{0}\\
     \mathbf{0} & A(G/\Pi)
     \end{bmatrix}=QB.
\end{equation*}
Conversely, suppose $AQ=QB$. %, then $Q^TAQ=B$ since $Q^TQ=I$.
Let $\e_a$ be the characteristic vector of vertex $a$ in $G$, and let $\check{\e}_a$ be the characteristic vector of vertex $a$ such that $Q\check{\e}_a=\frac{1}{\sqrt{2}}(\e_a-\e_{\hat{f}(a)})$ if $a\in V(X_1)$, and $Q\check{\e}_a$ is the normalized characteristic vector of the cell containing the vertex $a$, otherwise. If $a\in V(G)\backslash V(X_1)$ belong to $V_j$, then $\e_a^TQ=\e_b^TQ=\frac{1}{\sqrt{|V_j|}}\check{\e}_j^T$. Thus, if $k\in V(G)\backslash V(X_1)$, then 
$\e_a^T(AQ)\check{\e}_k=\e_a^T(QB)\check{\e}_k=\frac{1}{\sqrt{|V_j|}}\check{\e}_j^TB\check{\e}_k=\frac{1}{\sqrt{|V_k|}}\displaystyle\sum_{v\in V_k}A_{a,v}$. Hence,
\begin{center}
$\displaystyle\sum_{v\in V_k}A_{a,v}= \sqrt\frac{|V_k|}{|V_j|}\check{\e}_j^TB\check{\e}_k=c_{j,k}.$
\end{center}
As the above sum is independent of the choice of vertex in $V_j$, $\Pi$ is equitable. %the sum of edge weights from each vertex $a\in V_j$ to vertices in $V_k$ is a constant $c_{j,k}$ for all vertices in $V_j.$ Hence $\Pi$ is equitable. 

%\sarojini{Since $\e_a^T(AQ)\check{\e}_k=\frac{1}{\sqrt{|V_k|}}\displaystyle\sum_{v\in V_k}A_{a,v},$ it follows that $\e_a^T(AQ)\check{\e}_k=\e_a^T(QB)\check{\e}_k=\frac{1}{\sqrt{|V_j|}}\check{\e}_j^TB\check{\e}_k=\frac{1}{\sqrt{|V_k|}}\displaystyle\sum_{v\in V_k}A_{a,v}.$ Hence, $\displaystyle\sum_{v\in V_k}A_{a,v}= \sqrt\frac{|V_k|}{|V_j|}\check{\e}_j^TB\check{\e}_k=c_{j,k}$. So, it seems that $\frac{1}{\sqrt{|V_j|}}\check{\e}_j^TB\check{\e}_k$ should not equal to $\frac{1}{\sqrt{|V_j|}}c_{j,k}.$ }

Next, we prove $2\Leftrightarrow 3$. If 3 holds, then $AQ\check{\e}_u=QB\check{\e}_u$ for all $u$. Since $Q\check{\e}_u$ and $QB\check{\e}_u$ belong to the column space of $Q$, we obtain 2. The converse is straightforward.

Lastly, we prove $3\Leftrightarrow 4$.
If $AQ=QB$, then $Q^TA=BQ^T$ since $A$ and $B$ are symmetric. So, $AQQ^T=QBQ^T=QQ^TA$. Conversely, if $AQQ^T=QQ^TA$, then the fact $Q^TQ=I$ yields $Q^TAQQ^T=Q^TA$. If $B:=Q^TAQ$, then $BQ^T=Q^TA$, and because $A$ and $B$ are symmetric, we get $AQ=QB$. Finally, %if $F=\frac{1}{2}\begin{bmatrix}1&-1\\-1&1\end{bmatrix}$, then $F+\frac{1}{2}J_2=I_2$, and 
using the fact that $\mathcal{C}\mathcal{C}^T=\displaystyle\bigoplus_{k}\frac{1}{|V_k|}J_{|V_k|}$ gives us
\begin{equation*}
QQ^T%=[\mathcal{B}~~\mathcal{C}]\begin{bmatrix}\mathcal{B}^T\\\mathcal{C}^T\end{bmatrix}
=\mathcal{B}\mathcal{B}^T+\mathcal{C}\mathcal{C}^T=\left(\frac{1}{2}\bigoplus_{|V(X_1)|}\begin{bmatrix}1&-1\\-1&1\end{bmatrix}\oplus O\right)+\left(\bigoplus_{|V(X_1)|}\frac{1}{2}J_2\oplus \bigoplus_{k>|V(X_1)|} \frac{1}{|V_k|}J_{|V_k|}\right),
\end{equation*}
where $O$ is the zero matrix of the appropriate size. This yields the form of $QQ^T$ in 4.

Combining the above cases proves the equivalence of 1-4. Finally, if a matrix $B$ in 3 exists, then $\Pi$ is equitable. As $\mathcal{B}^T\mathcal{B}=I,$ we get $\mathcal{B}^TA\mathcal{B}=A(X_1)-A'$ and $\mathcal{C}^TA\mathcal{C}=A(G/\Pi).$ Now, if $a\in V(X_1)$ and $b\in V(G)\backslash V(X_1)$, then $Q\check{\e}_a= \frac{1}{\sqrt{2}}\ob{\e_a-\e_{\hat{f}(a)}}$ and $Q\check{\e}_b$ is a characteristic vector of $\Pi$. As $X_1$ and $X_2$ are edge-perturbed twin subgraphs in $G$,
\begin{center}
$\check{\e}_a^T(Q^TAQ)\check{\e}_b=(Q\check{\e}_a)^TA(Q\check{\e}_b)=\frac{1}{\sqrt{2}}\ob{\e_a-\e_{\hat{f}(a)}}^TA(Q\check{\e}_b)=0.$
\end{center} 
Therefore, $B=Q^TAQ=\begin{bmatrix} \mathcal{B}^TA\mathcal{B} & \mathcal{B}^TA\mathcal{C}\\ \mathcal{C}^TA\mathcal{B} & \mathcal{C}^TA\mathcal{C}\end{bmatrix}$, and because $\mathcal{B}^TA\mathcal{B}=A(X_1)-A'$ and $\mathcal{C}^TA\mathcal{C}=A(G/\Pi)$, we obtain the form of $B.$
%If $j,k\in V(G)\backslash V(X_1)$, then
%\begin{equation*}
%\begin{split}
%\check{\e}_j^T(Q^TAQ)\check{\e}_k&=(Q\check{\e}_j)^TA(Q\check{\e}_k)=\frac{1}{\sqrt{|V_j||V_k|}}\sum_{a\in V_j,b\in V_k}A_{a,b}=\frac{|V_j|}{\sqrt{|V_j||V_k|}}c_{j,k}=\sqrt{c_{j,k}c_{k,j}}.
%\end{split}
%\end{equation*}
%This shows that $A(G/\Pi)$ is the submatrix of $B$ indexed by $V(G)\backslash V(X_1)$. Now, if $a,b\in V(X_1)$, then the fact that $X_1$ and $X_2$ are edge-perturbed twin subgraphs in $G$ yields
%\begin{center}
%$\check{\e}_a^T(Q^TAQ)\check{\e}_b=(Q\check{\e}_a)^TA(Q\check{\e}_b)=\frac{1}{2}\ob{\e_a-\e_{\hat{f}(a)}}^TA\ob{\e_b-\e_{\hat{f}(b)}}=A_{a,b}-A_{a,\hat{f}(b)}.$
%\end{center}
%Lastly, if $a\in V(X_1)$ and $b\in V(G)\backslash V(X_1)$, then $Q\check{\e}_a= \frac{1}{\sqrt{2}}\ob{\e_a-\e_{\hat{f}(a)}}$ and $Q\check{\e}_b$ is a characteristic vector of $\Pi$. As $X_1$ and $X_2$ are edge-perturbed twin subgraphs in $G$,
%\begin{center}
%$\check{\e}_a^T(Q^TAQ)\check{\e}_b=(Q\check{\e}_a)^TA(Q\check{\e}_b)=\frac{1}{\sqrt{2}}\ob{\e_a-\e_{\hat{f}(a)}}^TA(Q\check{\e}_b)=0.$
%\end{center}
%Combining these observations yields the form of $B$. 
\end{proof}

\end{document}
